\chapter{Experiments}
\label{chap:experiments}

\section{Experimental Setup}
\label{sec:experimental-setup}

Experiments are organised around a staged comparison between heuristic
baselines and learned models for key-value pair extraction from business
documents. The overall goal is to evaluate not only prediction quality but
also the practicality of running the full pipeline from raw PDF sources to
benchmarked outputs.

For the learned baselines, each experiment is split into two operational
phases. First, raw KVP10k data are converted into page-level prepared samples
containing prompt text (Stage~3) or image-text-layout tensors (Stage~4),
target annotations, and aligned ground-truth key-value pairs. Second, the
model is trained on the prepared train split and evaluated on the prepared
test split.

% ---------------------------------------------------------------------------
\subsection{Hardware}
\label{sec:hardware}

\paragraph{Stage 1--3 (Mistral QLoRA baseline).}
Development and debugging were performed on CPU and login-node resources. The
final learned-baseline training runs were executed on the TinyGPU cluster
using a single NVIDIA A100~40\,GB GPU.

An important infrastructure constraint is that TinyGPU compute nodes do not
have internet access. Consequently, the Stage~3 workflow had to be split
operationally: PDF downloading and data preparation were executed on the login
node, while model training and prediction were executed offline on the GPU
node using a pre-populated HuggingFace cache.

The first corrected GPU run exposed a memory constraint on the available
hardware. A full bf16 LoRA setup exceeded the 40\,GB memory budget, so the
final training configuration was adapted to QLoRA with gradient checkpointing
and an 8-bit paged optimiser.

\paragraph{Stage 4 (LayoutLMv3\,+\,Linker).}
Stage~4 training is executed on the LRZ AI Systems cluster using A100~80\,GB
GPUs. Three separate SLURM jobs handle the Stage~4 workload:

\begin{enumerate}
  \item \textbf{Job 1 -- Data preparation:} CPU-only login node, converts
  KVP10k HuggingFace rows into Stage~4 image-tensor-annotation bundles,
  renders and caches page images at 150\,DPI.
  \item \textbf{Job 2 -- Stage~4a (entity pre-training):} Single A100~80\,GB
  GPU, 24-hour wall-clock limit, trains LayoutLMv3 encoder + entity head
  with the linker disabled.
  \item \textbf{Job 3 -- Stage~4b (joint fine-tuning):} Single A100~80\,GB
  GPU, 12-hour wall-clock limit, resumes from Stage~4a checkpoint and
  activates the biaffine linker with $\lambda > 0$.
\end{enumerate}

All Stage~4 computations use bfloat16 mixed precision, requiring no
quantisation of the 125M-parameter LayoutLMv3-base model. Peak GPU memory
utilisation during training is approximately 42\,GB (encoder forward pass +
gradient accumulation buffer), well within the 80\,GB envelope.

% ---------------------------------------------------------------------------
\subsection{Hyperparameters}
\label{sec:hyperparameters}

\subsubsection{Stage~3: Mistral QLoRA Baseline}
\label{sec:stage3-hyperparams}

The learned baseline follows the corrected IBM-style configuration as closely
as possible while remaining feasible on the available hardware.
The main settings are summarised in Table~\ref{tab:mistral-hyperparams}.

\begin{table}[h]
  \centering
  \caption{Main hyperparameters for the corrected Mistral baseline.}
  \label{tab:mistral-hyperparams}
  \begin{tabular}{ll}
    \hline
    \textbf{Parameter} & \textbf{Value} \\
    \hline
    Base model & Mistral-7B-Instruct-v0.2 \\
    Adaptation method & LoRA on 4-bit QLoRA base model \\
    LoRA rank $r$ & 4 \\
    LoRA $\alpha$ & 4 \\
    LoRA dropout & 0.05 \\
    Learning rate & $5 \times 10^{-4}$ \\
    Epochs & 8 \\
    Per-device batch size & 1 \\
    Gradient accumulation & 4 \\
    Maximum sequence length & 8192 \\
    Optimiser & paged\_adamw\_8bit \\
    Precision / compute dtype & bfloat16 compute on quantised weights \\
    Gradient checkpointing & enabled \\
    Decoding & greedy \\
    \hline
  \end{tabular}
\end{table}

\paragraph{QLoRA and LoRA rank justification.}
The original full-precision LoRA configuration exceeded the 40\,GB memory
budget of the available A100 GPU. To make the experiment feasible on available
hardware, we adopted QLoRA with 4-bit NF4 quantisation (rank $r = 4$,
alpha $\alpha = 4$). While rank 4 is lower than the estimated $r \geq 16$ in
the official KVP10k baseline, this is a necessary hardware adaptation. The
QLoRA paper (Dettmers et al., 2023) demonstrates that 4-bit quantisation does
not degrade model performance relative to 16-bit LoRA on standard benchmarks,
so quantisation is not the primary limitation. However, the reduced rank and
55.8\% training data coverage (581 of 5{,}389 prepared unique pages, due to
PDF download and text extraction failures) both reduce the model's adaptation
capacity and may contribute to performance differences relative to the
published baseline.

\subsubsection{Stage~4: LayoutLMv3\,+\,Linker}
\label{sec:stage4-hyperparams}

Table~\ref{tab:stage4-hyperparams} summarises the Stage~4 hyperparameters.

\begin{table}[h]
  \centering
  \caption{Hyperparameters for Stage~4 (LayoutLMv3\,+\,Biaffine Linker).}
  \label{tab:stage4-hyperparams}
  \begin{tabular}{ll}
    \hline
    \textbf{Parameter} & \textbf{Value} \\
    \hline
    Base model & microsoft/layoutlmv3-base \\
    Encoder parameters & 125M (all unfrozen) \\
    Entity head & 2-layer MLP, hidden dim 768, GELU \\
    Linker head & Biaffine, weight matrix $768 \times 768$ \\
    Linking weight $\lambda$ & swept over \{0.5, 1.0, 2.0\} \\
    Learning rate (Stage~4a) & $2 \times 10^{-5}$ \\
    Learning rate (Stage~4b encoder) & $2 \times 10^{-6}$ \\
    Learning rate (Stage~4b linker) & $2 \times 10^{-5}$ \\
    Per-device batch size & 4 \\
    Gradient accumulation steps & 8 \\
    Effective batch size & 32 \\
    Warmup steps & 500 \\
    Maximum epochs & 20 \\
    Early stopping patience & 3 \\
    Optimiser & AdamW, weight decay 0.01 \\
    Precision & bfloat16 (no quantisation) \\
    Max sequence length & 512 tokens \\
    Image resolution & $224 \times 224$ \\
    Gradient checkpointing & enabled \\
    \hline
  \end{tabular}
\end{table}

\paragraph{$\lambda$ sweep rationale.}
The linking loss weight $\lambda$ controls the relative influence of the
linker objective on the shared encoder representations. A value that is too
large risks distorting the entity classification capacity; too small and the
linker receives insufficient gradient signal. We sweep $\lambda \in
\{0.5, 1.0, 2.0\}$ on a held-out 10\% validation split of the training set
(Stage~4b only) and select the value that maximises combined F1.

% ---------------------------------------------------------------------------
\section{Baseline Models}
\label{sec:baseline-models}

\subsection{Nearest Neighbour Baseline}
\label{sec:nn-baseline}

Our heuristic baseline uses a spatial proximity heuristic to pair keys with
values:
\begin{enumerate}
  \item Extract all annotations with linking attributes (keys).
  \item Extract all annotations with text content (potential values).
  \item For each key, find the nearest value within maximum distance
  $d_{\max} = 0.3$ (Euclidean distance between bounding box centroids,
  normalised by image dimensions).
  \item If a value is found, create a KV pair; otherwise mark as
  ``unvalued''.
\end{enumerate}

\subsection{Baseline Performance}
\label{sec:baseline-performance}

We evaluate the baseline on a 20-sample subset for initial validation:

\begin{table}[h]
  \centering
  \caption{Nearest Neighbour Baseline Results (20-sample validation).}
  \label{tab:nn-baseline}
  \begin{tabular}{llll}
    \hline
    \textbf{Metric} & \textbf{Precision} & \textbf{Recall} & \textbf{F1} \\
    \hline
    Combined & 0.093 & 1.000 & 0.169 \\
    Text-only & 0.093 & 1.000 & 0.169 \\
    Location-only & 0.093 & 1.000 & 0.169 \\
    \hline
  \end{tabular}
\end{table}

High recall (100\%) reflects the fact that the baseline pairs all detected
keys; low precision (9.3\%) arises from many false positive pairings.
$F_1 = 0.169$ establishes the lower bound for future models.

% ---------------------------------------------------------------------------
\section{Main Benchmark Evaluation}
\label{sec:main-benchmark}

\subsection{Test Set Results}
\label{sec:test-set-results}

The Stage~4 model is evaluated on the same 581-page prepared test set used
for the Stage~3 Mistral baseline (Section~\ref{sec:dataset-coverage}), enabling
direct comparison. Evaluation follows the dual-threshold protocol:
$\text{NED} < 0.2$, $\text{IoU} > 0.3$, combined F1 as the primary metric.

Table~\ref{tab:stage4-results-placeholder} reports Stage~4 results once
experiments are complete.

\begin{table}[h]
  \centering
  \caption{Stage~4 LayoutLMv3\,+\,Linker Results
    (NED\,=\,0.2, IoU\,=\,0.3; $n = 581$ test pages).}
  \label{tab:stage4-results-placeholder}
  \begin{tabular}{lllllll}
    \hline
    \textbf{Mode} & \textbf{Prec.} & \textbf{Rec.} & \textbf{F1}
                  & \textbf{TP} & \textbf{FP} & \textbf{FN} \\
    \hline
    Text-only (NED) & -- & -- & -- & -- & -- & -- \\
    Text+Bbox (NED+IoU) & -- & -- & -- & -- & -- & -- \\
    Linking F1 & \multicolumn{6}{c}{-- (key-value association only)} \\
    \hline
  \end{tabular}
\end{table}

\subsection{Per-Cluster Analysis}
\label{sec:per-cluster-analysis}

Following the per-cluster breakdown established for Stage~3
(Section~\ref{sec:stage3-per-cluster}), Stage~4 results are reported
separately for Cluster~0 (dense layouts, $n = 465$ documents) and Cluster~1
(sparse layouts, $n = 116$ documents). This analysis directly tests whether
LayoutLMv3's spatial grounding addresses the sparse-layout hallucination
problem identified in Stage~3.

\begin{table}[h]
  \centering
  \caption{Stage~4 Per-Cluster Results (to be filled after experiments).}
  \label{tab:stage4-cluster-results}
  \begin{tabular}{llllll}
    \hline
    \textbf{Cluster} & \textbf{Docs} & \textbf{Pred} & \textbf{GT}
                     & \textbf{Correct (\%)} & \textbf{Halluc. (\%)} \\
    \hline
    Cluster~0 (Dense)  & 465 & -- & -- & -- & -- \\
    Cluster~1 (Sparse) & 116 & -- & -- & -- & -- \\
    \hline
  \end{tabular}
\end{table}

\subsection{Error Analysis}
\label{sec:stage4-error-analysis}

The error analysis for Stage~4 follows the same taxonomy applied to Stage~3:
hallucination rate, missed-entity rate, and spatial grounding errors.
In addition, linking errors are classified into three types:
\begin{itemize}
  \item \textbf{False-positive link:} a predicted link whose key or value
  entity does not match any ground-truth entity.
  \item \textbf{False-negative link:} a ground-truth link not predicted by
  the model.
  \item \textbf{Mis-link:} a correctly predicted key entity paired with
  an incorrect value entity.
\end{itemize}

% ---------------------------------------------------------------------------
\section{Robustness Experiments}
\label{sec:robustness-experiments}

\subsection{Rotation Robustness}
\label{sec:rotation-robustness}

We evaluate robustness to document rotation at angles:
$\{0^\circ, 5^\circ, 10^\circ, 15^\circ\}$. Rotations are applied to the
document image before rendering; bounding box coordinates are updated via
the corresponding 2D rotation transform.

\subsection{Occlusion Robustness}
\label{sec:occlusion-robustness}

We simulate occlusion by masking random rectangular regions covering
$\{0\%, 10\%, 20\%, 30\%\}$ of the document area. Masked regions are
filled with the mean image colour to avoid introducing distribution-shift
artifacts beyond the intended occlusion.

\subsection{Annotation Dropout}
\label{sec:annotation-dropout}

We evaluate with annotation dropout rates of
$\{0\%, 10\%, 20\%, 30\%\}$: randomly selected ground-truth annotations
are removed from the supervision signal at training time, simulating
incomplete labelling. At test time, all annotations are retained.

% ---------------------------------------------------------------------------
\section{Ablation Studies}
\label{sec:ablation-studies}

\subsection{Linker Module}
\label{sec:ablation-linker}

We train a model variant with the biaffine linker disabled ($\lambda = 0$)
to isolate the contribution of the linking objective. In this variant, key-value
association falls back to a greedy nearest-neighbour heuristic applied to the
classified entity spans, as a direct like-for-like comparison with Stage~3.

\subsection{Spatial Encoding}
\label{sec:ablation-spatial}

We ablate the bounding-box spatial augmentation in the linker
(Section~\ref{sec:spatial-encoding}) by setting $\mathbf{W}_{\text{sp}} = 0$,
forcing the linker to rely solely on the encoder's learnt spatial representations.
This tests whether the explicit injection of coordinate features provides
additional benefit beyond what LayoutLMv3's layout embeddings already capture.

\subsection{Pre-training}
\label{sec:ablation-pretraining}

We compare the full pre-trained LayoutLMv3-base initialisation against a
random initialisation of the encoder weights (keeping the tokenizer and
bounding-box embedding tables fixed). This ablation quantifies the benefit
of the multimodal pre-training described in
Section~\ref{sec:pretraining-objectives}.
